% Abstract page header, per WSU guidelines. Title copied exactly from the title page, no bold or underline.
\begin{center}
  {\doublespacing\MakeUppercase{\thesistitle}\par}
  \vspace{3\baselineskip}
  Abstract\par
  \vspace{\baselineskip}
  {\singlespacing
  \thesisauthor, Ph.D.\\
  Washington State University\\
  \thesisdate\par}
\end{center}
\vspace{\baselineskip}
\noindent Chair: Brian A. Collins
\addcontentsline{toc}{chapter}{Abstract}

\vspace{\baselineskip}
\noindent The emergent electrical optical and kinetic properties of many materials are set by the interfaces between buried structures and assemblies hidden within the material. Atoms and molecules there can arrange and bond differently than in the bulk, yet these layers are inaccessible and often only one or two molecules thick. Measuring them requires sensitivity to both structure and electronic state, with sub-nanometer resolution. Resonant soft X-ray scattering (RSoXS) is sensitive to both. Its contrast is set by specific chemical bonds and, with polarized X-rays, by their orientation. In reflection geometry, this orientation can be resolved with depth. In practice, the anisotropic optical tensor is fit freely at each energy and in each layer, yet this analysis has many parameters and only a loose connection to molecular structure. Here we build the tensor from the electronic transitions of a single molecule and use it to recover the molecular state at buried interfaces.

Transitions from density functional theory are clustered and refined against angle-resolved NEXAFS (ML-DFT). The tensor of a film is a Kramers-Kronig consistent sum over these transitions, weighted by molecular orientation and density. With this tensor, polarized resonant reflectivity of vapor-deposited zinc phthalocyanine (ZnPc) is fit as well as with a free tensor, using an order of magnitude fewer parameters. No contamination or oxidation layer is needed. Within one molecular layer of the native oxide, the molecules lie face-on at reduced density. At the free surface they are denser and tilted near the magic angle. We attribute this profile to competition between substrate adhesion and \(\pi - \pi\) cohesion during growth. If so, a partial film measured at a coverage checkpoint is in a transient state.

Refining the amplitude and dipole angle of each resonance with depth gives a spectrum for each layer of the film. Films of tris(8-hydroxyquinolinato)aluminum (Alq\textsubscript{3}) made under nominally identical conditions have different NEXAFS spectra. Analyzing one film with the spectrum of another gives an unphysical substrate density, and refining the amplitudes with depth removes this error. The electronic state of a molecule in a film cannot be taken from a reference spectrum of another film.

Reference NEXAFS spectra are rarely shared with a record of how they were measured or processed. With collaborators at seven institutions, we built X-ray Atlas, an open, programmatic database of NEXAFS spectra and optical models. Every contribution is normalized and scored the same way, and every processing revision has a DOI. Each optical model is tied to a specific data revision, so a tensor such as the one for ZnPc can be cited, reproduced, and extended.

Molecular orientation and electronic state can now be measured at the buried interfaces of finished films. The same model can be built for any molecule whose transitions can be calculated. With shared, versioned optical models, these measurements can be compared between samples and between laboratories.
\clearpage
